\subsubsection{库存稳定性与组织均衡的共同解释}
当前两组选择将S006单独执行，其余14区构成另一组，工作量变异系数为0.941762。该分组在资源需求上逐项不高于另两种，总配置也从33降为28，所以按“缺口—总配置—均衡”的既定顺序具有支配依据。库存改变时的稳定性来自函数$\sum_p(D_p-I_p)_+$对$D$的分量单调性，不只是有限试验中恰好没有切换。

均衡指标仍有实际意义：最节省装备的方案把较多任务集中在一组，需要相应协调能力。若把均衡置于资源投入之前，决策偏好已经改变，应重新比较三种合法分区。把选择依据与工作量代价同时列出，能让评阅者看懂“为什么如此分组”，而不必把任何一个分组宣传为所有条件下都占优。

\subsubsection{上游预案与独立配置的衔接}
问题四内部固定日程后可以完整枚举；改变上游预案则需重新构造占用区间和通信关联。计算顺序是：验证新联合预案，冻结它，构建不可拆块，再核算各组资源。只替换总工期而沿用旧资源峰值，会丢失任务重叠和保障关系的变化。

面向三组的折中预案把三组缺口从7降到6，代价是工期增加29.171 s，两组缺口同时从1升到2。较短工期、较少两组资源和较少三组资源因此不是同一个目标。执行前已经明确组织方式时，可利用这一候选比较选择对应预案；尚未确定时，则保留多个已验证预案及其资源清单。

进一步把分组配置代价反馈给上游位置和时间决策，将形成更强的运输—通信—组织联合设计。当前工作已经提供完整接口和有限候选上的量化比较，为这类扩展确定了可调整对象与核验条件。
